\documentclass[12pt, a4paper]{article}
\usepackage[T1]{fontenc}
\usepackage{lmodern}
\usepackage{amsmath, amsthm, amssymb, mathtools}
\usepackage[margin=1in]{geometry}
\usepackage{xr-hyper}
\usepackage[colorlinks=true,linkcolor=blue,citecolor=blue]{hyperref}
\usepackage{cleveref}
\usepackage{graphicx, subcaption}
\usepackage{natbib}
\usepackage{booktabs}
\usepackage{enumerate}

% Code references: scaled monospace with line-breaking
\newcommand{\code}[1]{\texttt{\small #1}}

\newtheorem{theorem}{Theorem}[section]
\newtheorem{lemma}[theorem]{Lemma}
\newtheorem{proposition}[theorem]{Proposition}
\newtheorem{corollary}[theorem]{Corollary}
\newtheorem{definition}{Definition}[section]
\newtheorem{remark}{Remark}[section]
\newtheorem{conjecture}[theorem]{Conjecture}
\newtheorem{derivation}[theorem]{Derivation}
\newtheorem{observation}[theorem]{Observation}
\numberwithin{equation}{section}

\newcommand{\dd}{\mathrm{d}}
\newcommand{\seriestitle}{Why Rotating Fluids Make Polygons}


\externaldocument[I-]{../paper-1-mathematics/main}
\externaldocument[II-]{../paper-2-physics/main}
\externaldocument[III-]{../paper-3-gravity/main}
\externaldocument[IV-]{../paper-4-field-theory/main}
\externaldocument[V-]{../paper-5-cosmology/main}

\title{\seriestitle\\[6pt]\large Supplement S: Extended Proofs and Computations}
\author{Gordon Speagle}
\date{}

\begin{document}
\maketitle

\begin{abstract}
This supplement collects extended proofs, detailed computations,
and numerical convergence analyses that support the results of
Papers~I--V.
Each section is self-contained and referenced from the
corresponding main paper.
\end{abstract}

%% ====================================================================
\section{Paper~I supplements: K-theory details}
\label{sec:kasparov}

%% Content will be moved here from Paper I appendix

\subsection{The Kasparov product construction}
\label{sec:kasparov-construction}

%% Placeholder: Kasparov product proof from Paper I

\subsection{Bolza surface convergence analysis}
\label{sec:bolza-convergence}

%% Placeholder: Bolza convergence narrative from Paper I

%% ====================================================================
\section{Paper~II supplements: multi-ring proof details}
\label{sec:multiring}

%% Placeholder: multi-ring casework from Paper II

%% ====================================================================
\section{Paper~III supplements: BTZ spectral statistics}
\label{sec:btz-lorentzian}

\paragraph{Lorentzian continuation and the BTZ horizon.}
The Euclidean analysis above determines the thermal phase structure; the Lorentzian continuation identifies the real-time dynamics.
The Green's function offset $\Phi(\rho) = \log(2\sinh\rho) + b(N)$ admits
an analytic continuation $\rho \to \rho + i\pi/2$ that interchanges
the BTZ exterior and interior:
$\sinh(\rho + i\pi/2) = i\cosh\rho$, so
$\Phi^{\mathrm{int}} = \log(2\cosh\rho) + b(N) + i\pi/2$.
The analytically continued kernel interpolates between curvature
signs:
\begin{center}
\begin{tabular}{lccc}
\toprule
$\mathrm{Im}(\rho)$ & $0$ & $\pi/2$ & $\pi$ \\
\midrule
Curvature sign & $K{=}{-}1$ & $K{=}0$ & $K{=}{+}1$ \\
Continued kernel & $\log(2\sinh\rho)$ & $\log(2\cosh\rho) + i\pi/2$ &
  $\log(2\sinh\rho) + i\pi$ \\
\bottomrule
\end{tabular}
\end{center}
\noindent
(The real parts of the continued kernels reproduce the
pair interactions on spaces of the corresponding curvature
sign; the imaginary parts are phases.)
At each palindromic threshold $\rho^*(N)$: the critical eigenvalue
$\lambda_{m^*}$ crosses zero (\emph{real} side), while the interior
resonance has $\mathrm{Im}(\lambda) = \pi/2$ (universal, independent
of $N$ and~$m$).
The polygon--BTZ transition is a \emph{time--space exchange}:
for $\rho > \rho^*$ the mode oscillates in time
($\omega = \sqrt{\lambda_{m^*}}$); for $\rho < \rho^*$ it grows
spatially ($\kappa = \sqrt{|\lambda_{m^*}|}$);
at $\rho = \rho^*$ the Killing vector degenerates.

Under Lorentzian continuation ($\rho \to \rho + i\varepsilon$),
the partition function acquires a branch cut at the threshold:
$Z_L \sim A(N)\,({\rho - \rho^* + i\varepsilon})^{-1/2}$
with spectral function $\mathrm{Im}[Z_L] \propto |\rho - \rho^*|^{-1/2}$
(a van~Hove singularity).
The exponent $-1/2$ (divergent density at the edge) is the
spectral signature of the integrable Havelock system.
The Dyson circular ensemble (GUE, $\beta = 2$) has the
opposite exponent: $\rho \propto |\lambda - \lambda_{\mathrm{edge}}|^{+1/2}$
(vanishing density, the Wigner semicircle edge).
The sign flip $-1/2 \to +1/2$ is a theorem:
the Havelock gaps satisfy $\lambda_{m^*+\delta} = \delta^2/2$
(even~$N$) or $\delta(\delta{-}1)/2$ (odd~$N$)
(verified numerically for
$N = 7,\ldots,30$; see Code Availability), giving $\rho(\lambda) = (2\lambda)^{-1/2}$
(van~Hove, the integrable exponent \citep{BerryTabor1977}).
The GUE edge exponent $+1/2$ is the Tracy--Widom theorem
\citep{TracyWidom1994}.
The BTZ black hole has GUE level statistics
\citep{AurichSteiner1993}.


%% ====================================================================
\section{The Wheeler--DeWitt equation on the palindromic minisuperspace}

%% ====================================================================
\section{Paper~V supplements: Pell identity proof}
\label{sec:pell-proof}

\section{Proof of the Pell identity}
\label{app:pell}

\begin{proof}[Proof of the Pell identity~\eqref{eq:pell-identity}
  via the Klein quartic]
The argument has four steps.

\emph{Step~1 (The Pell geodesic on the modular surface).}
The Pell equation $8^2 - 7 \times 3^2 = 1$ defines the matrix
$M = \bigl(\begin{smallmatrix} 8 & 21 \\ 3 & 8 \end{smallmatrix}\bigr)
\in \mathrm{SL}(2,\mathbb{Z})$
with $\det M = 1$, $\operatorname{tr} M = 16$.
This matrix is hyperbolic and defines a closed geodesic on the
modular surface $\mathrm{SL}(2,\mathbb{Z})\backslash\mathbf{H}^2$
of length
$\ell = 2\operatorname{arccosh}(\operatorname{tr}M/2)
= 2\operatorname{arccosh}(8) = 2\ln\varepsilon_7$.

\emph{Step~2 (Lift to the Klein quartic).}
The level-$7$ modular curve
$X(7) = \Gamma(7)\backslash\mathbf{H}^2$
is the \emph{Klein quartic}---a genus-$3$ Riemann surface with
automorphism group $\mathrm{PSL}(2,\mathbb{F}_7)$, order~$168$.
This group contains $\mathbb{Z}/7\mathbb{Z}$ as a subgroup,
and $X(7)/(\mathbb{Z}/7\mathbb{Z})
= \mathrm{SL}(2,\mathbb{Z})\backslash\mathbf{H}^2$.
The Pell geodesic on the modular surface lifts to
\emph{seven} closed geodesics on $X(7)$, one per
$\mathbb{Z}_7$~coset, each of length~$2\ln\varepsilon_7$.

\emph{Step~3 (Selberg trace formula on $X(7)$).}
The BO potential $V_7(\rho) = \ln(2\sinh\rho) - u^*$ is the
radial effective potential of the Laplacian on~$\mathbf{H}^2$
in the $\mathbb{Z}_7$-equivariant sector.
The WKB spectral sum of this potential equals the
spectral side of the Selberg trace formula on~$X(7)$,
evaluated with the WKB test function~$h$.
On the \emph{geometric} side, each primitive geodesic of
length~$\ell$ contributes
$\ell/(2\sinh(\ell/2)) \times \hat h(\ell)$
to the trace.
For the Pell geodesic ($\ell_0 = 2\ln\varepsilon_7$):
the WKB spectral factor is
$\hat h(\ell_0) = 2\sinh(\ell_0/2) = 2\sinh(\operatorname{arccosh}8) = 6\sqrt{7}$,
which \emph{exactly cancels} the Selberg stability denominator
$1/(2\sinh(\ell_0/2))$.
We prove this cancellation explicitly.
The potential $V_7(\rho) = \ln(2\sinh\rho)$ is (up to a constant)
the radial Green's function of the Laplacian on~$\mathbf{H}^2$.
Its Legendre spectral coefficient is the resolvent
\citep[Ch.~IV]{Hejhal1976}:
\[
  \int_0^\infty \ln(2\sinh\rho)\,
  P_{-1/2+ir}(\cosh\rho)\,\sinh\rho\,d\rho
  = \frac{-\pi}{2r\sinh(\pi r)}.
\]
The WKB test function for the Selberg trace formula is
$h(r) = (1/4 + r^2) \times a(r)$, where $a(r)$ is the
spectral coefficient and $(1/4 + r^2)$ is the Laplacian eigenvalue.
The Selberg transform $\hat{h}(\ell)$ is the Abel--Fourier
transform of~$h$; for $h(r) \propto (1/4 + r^2)/(r\sinh(\pi r))$,
this evaluates via the residue formula
$\sum \operatorname{Res}[\hat{h}]$ at the poles $r = in$
($n = 0, 1, 2, \ldots$) of $1/\sinh(\pi r)$:
\[
  \hat{h}(\ell) = \frac{1}{2\pi}\int_{-\infty}^{\infty}
  e^{ir\ell}\,\frac{\pi(1/4+r^2)}{2r\sinh(\pi r)}\,dr
  = 2\sinh(\ell/2).
\]
The last equality follows from the Mittag-Leffler expansion
of $1/\sinh(\pi r)$ and the factorisation
$(1/4 + r^2) = (r + i/2)(r - i/2)$, which shifts the
poles of $1/(r\sinh(\pi r))$ to give the exponentials
$e^{\pm\ell/2}$ whose difference is $2\sinh(\ell/2)$.
This completes the proof: $\hat{h}(\ell) = 2\sinh(\ell/2)$
cancels the stability denominator $1/(2\sinh(\ell/2))$,
leaving a net contribution of~$\ell_0$ per geodesic orbit.
The net contribution per orbit is therefore simply
the \emph{geodesic length}~$\ell_0 = 2\ln\varepsilon_7$.
Summing the seven lifts:
\[
  2S_{\mathrm{BO}}(7)
  = 7 \times \ell_0 - \text{quantum corrections}
  = 14\,\ln\varepsilon_7 - \text{corrections}.
\]

\emph{Step~4 (Quantum corrections from the Dunham series).}
The Maslov index of the bounce orbit gives the correction
$-\sqrt{2/\pi}$ (the WDW zero-point energy).
The anharmonicity of $\ln(2\sinh\rho)$ near $\rho = 0$
gives the Dunham correction $-4/c_7^2$ (verified to $0.26\%$
at $50$-digit precision).
Together: $2S/\ln\varepsilon_7
= 14 - \sqrt{2/\pi} - 4/c_7^2 + O(1/c_7^3)$.
\end{proof}

%% ====================================================================

\bibliographystyle{plainnat}
\bibliography{../shared/refs}

\end{document}
